% Replicability appendix. Included from decoder_paper_v2.tex via \input.
% Follows the NeurIPS / Sensors-journal reproducibility-checklist conventions.

\section{Replicability Checklist}
\label{app:replicability}

This section documents implementation details necessary for an independent
team to reproduce the headline results. The checklist follows the NeurIPS
and \emph{Sensors} reproducibility-checklist conventions.

\subsection{Code and Data}
\label{app:repl-code}

\begin{itemize}
\item \textbf{Code repository.} The full source tree, including
preprocessing, training, evaluation, and figure-generation scripts, is
released alongside the manuscript. Critical scripts:
\texttt{scripts/evaluation/run\_decoder\_quick\_test.py} (training
driver), \texttt{scripts/analysis/aggregate\_v8\_results.py}
(metric aggregator), and the experiment shells under
\texttt{scripts/experiments/}.
\item \textbf{Dataset.} The 120-file V8 corpus is released under the same
license as the encoder paper~\cite{encoder_paper}. Per-fold preprocessed
NPZ files are available at \TBD{data\_release\_url}.
\item \textbf{Hyperparameter sweep records.} All 34 cells of the 3-stage
HP sweep are recorded under \texttt{outputs/decoder20260511/checkpoints/}
with per-cell metrics.json. The composite-winner selection logic is in
\texttt{scripts/analysis/aggregate\_hp\_sweep\_all\_stages.py}.
\end{itemize}

\subsection{Random Seeds and Determinism}
\label{app:repl-seeds}

\begin{itemize}
\item \textbf{Decoder training seed.} 42 unless noted (all reported
5-fold experiments).
\item \textbf{Seed-variance estimate.} Three additional seeds were
trained as part of Stage 3 of the HP sweep on the Stage-1 winner config:
\{123, 456, 2024\}. Command-accuracy spans 0.293--0.337 across the three
seeds, vs.\ the seed-42 baseline of 0.325. The reported headline number
falls within this seed-variance envelope.
\item \textbf{Data-shuffle seed.} The data loader uses
\texttt{torch.Generator.manual\_seed(seed)} where applicable. CUDNN
deterministic mode is NOT enforced (the small non-determinism it adds
is dominated by the seed-variance envelope reported above).
\end{itemize}

\subsection{Hardware}
\label{app:repl-hardware}

\begin{itemize}
\item \textbf{GPUs.} NVIDIA RTX A6000, 48~GB VRAM. Two GPUs were used
during the HP sweep with parallel dispatch (one decoder cell per GPU).
\item \textbf{CPU + RAM.} The host has approximately 100 CPU cores and
1~TB of system RAM. The encoder-memory caching step uses up to ~16~GB
RAM per cell; the full multi-cell parallel sweep peaked at ~120~GB.
\item \textbf{Storage.} The corpus, checkpoints, audit JSONs, and W\&B
logs occupy approximately 80~GB of disk.
\end{itemize}

\subsection{Training Time and Compute Budget}
\label{app:repl-compute}

\begin{itemize}
\item \textbf{Single fold, decoder only.} ~30 minutes on one A6000
at the headline configuration ($d_{\text{model}}=384$, $n_{\text{layers}}=8$,
batch size 64, max\_token\_len=32 or 1400 depending on target mode,
300 epochs cap, patience 75).
\item \textbf{Five-fold cross-validation, decoder only.} ~2.5~hours
sequential on one GPU.
\item \textbf{Hyperparameter sweep (34 cells).} ~12--15~hours wall-clock
on two GPUs in parallel.
\item \textbf{End-to-end remediation (this paper).} ~3 days of GPU
compute including the audit phase, two preprocessing modes, the HP
sweep, the headline 5-fold, ablations, and Phase F preparation.
\end{itemize}

\subsection{Software Versions}
\label{app:repl-software}

\begin{itemize}
\item Python 3.10
\item PyTorch 2.x with CUDA 12.x
\item scikit-learn (for HGB baseline + metric reporting + StandardScaler);
SciPy (for ANOVA); XGBoost (failed on CPU-only host, replaced with sklearn
HGB --- see Section~\ref{sec:results-baseline})
\item Weights and Biases (project \texttt{gcode-decoder-2026}, public link
\TBD{wandb\_project\_url}) for hyperparameter-sweep tracking
\end{itemize}

\subsection{Hyperparameter Ranges Explored}
\label{app:repl-hp-ranges}

The 3-stage HP sweep explored the following ranges; the composite winner
within each is bolded.

\begin{itemize}
\item Learning rate: \{$1\!\times\!10^{-4}$, $\mathbf{5\!\times\!10^{-5}}$, $2\!\times\!10^{-5}$, $1\!\times\!10^{-5}$, $5\!\times\!10^{-6}$\}
\item Batch size: \{$\mathbf{64}$, 128\}
\item Model dimension $d_{\text{model}}$: \{$\mathbf{384}$, 512\}
\item Decoder layers: \{$\mathbf{8}$, 10, 12\}
\item Dropout: \{0.05, $\mathbf{0.1}$, 0.2, 0.3\}
\item Weight decay: 0.05 (fixed)
\item Legacy token loss weight: \{1.0, $\mathbf{3.0}$\}
\item Digit loss weight: \{$\mathbf{1.0}$, 3.0\}
\item Warmup epochs: \{0, $\mathbf{10}$\}
\item Scheduled sampling probability: \{0.0, 0.1, 0.2, $\mathbf{0.5}$, 1.0\}
\item Label smoothing: \{0.0, 0.1, 0.2\} (winner: 0.0)
\item Focal-loss $\gamma$: \{0.0, 2.0\} (winner: 0.0)
\item Curriculum schedule: \{none, 3-phase\} (winner: none)
\item Output mode: \{per\_row, $\mathbf{full\_window}$\}
\item Encoder fine-tuning: \{frozen, e2e at $1\!\times\!10^{-5}$, e2e at $5\!\times\!10^{-6}$\} (winner: frozen; e2e diverged at lr$\ge\!1\!\times\!10^{-5}$)
\item Position metadata exposed: \{$\mathbf{false}$, true\}
\end{itemize}

\subsection{Statistical Methodology}
\label{app:repl-stats}

\begin{itemize}
\item \textbf{Cross-validation.} File-level 5-fold stratified by operation
class. G-code-line coverage between train and test splits is repaired via
the algorithm in Appendix~\ref{app:coverage-repair}.
\item \textbf{ANOVA.} One-way ANOVA across the 5 folds of each ablation
versus the no-shortcuts baseline. Implemented with
\texttt{scipy.stats.f\_oneway}.
\item \textbf{Effect size.} Cohen's $d$ using pooled variance.
\item \textbf{Bootstrap CI.} 10{,}000 percentile-method resamples on the
5-fold means.
\item \textbf{Multiple-comparisons correction.} Holm-Bonferroni and BH-FDR
applied across all (ablation, macro-metric) pairs reported in
Table~\ref{tab:anova}, AND separately across the drilled per-class /
per-axis / per-position grid below. Significance at $\alpha=0.05$ is
reported under both adjustments.

\item \textbf{Per-class / per-axis / per-position ANOVA grid.} In
addition to the macro-metric ANOVA above (token / sequence / type /
command / param-type / sign / numeric), we run the same one-way-ANOVA +
Cohen's $d$ + Holm-Bonferroni / BH-FDR procedure on the drilled
per-class metrics: precision, recall, and F1 for each command class
(G0 / G1 / G2 / G3 / G53 / M30), each param-type axis (X / Y / Z / R /
F), each sign class (POSITIVE / NEGATIVE), each type class (SPECIAL /
COMMAND / PARAM / NUMERIC), each of the six numeric-digit output
positions (accuracy + macro precision / recall / F1), and each of the
ten digit-value classes (0--9). The grid contains approximately
\TBD{n\_per\_class\_anova\_tests} (ablation, drilled-metric) tests;
results are emitted to \texttt{audit/anova\_per\_class\_results.json}.
Per-class results with $n_{\text{support}} \le 50$ are flagged as
underpowered (Section~\ref{sec:lim-power}); the headline-significance
claims rely on the larger-support classes only.
\end{itemize}

\subsection{Inference Cost}
\label{app:repl-inference}

\begin{itemize}
\item \textbf{Decoder forward pass.} $\sim$\TBD{decoder\_flops\_per\_token} FLOPs per generated token.
\item \textbf{Single-row inference latency.} $\sim$\TBD{decoder\_latency\_ms\_row} ms on one A6000 at batch size 1, sequence length 32 (per-row mode) or up to 1{,}400 (full-window mode).
\item \textbf{Memory footprint.} The decoder + frozen encoder uses approximately \TBD{decoder\_mem\_mb}~MB of VRAM at inference time.
\end{itemize}
